% TOP_PROBLEM: {"id": 30007211, "title": "QMA hardness of continuum electronic energy with binary nuclear charges", "category_id": 16, "category": {"id": 16, "name": "physics", "display_name": "Mathematical Physics", "description": "Problems at the intersection of mathematics and physics.", "slug": "physics", "order_index": 16, "created_at": "2026-07-31T15:26:25.671Z"}, "status": "open"}
% FIRST_POSTED: 2026-10-06
% ENTRY_KIND: statement_only
% !TeX program = lualatex
\documentclass[11pt]{article}
\usepackage[margin=25mm]{geometry}
\usepackage{fontspec}
\setmainfont{Latin Modern Roman}
\usepackage{amsmath,amssymb,mathtools}
\usepackage{unicode-math}
\setmathfont{Latin Modern Math}
\usepackage{xurl,hyperref}
\hypersetup{colorlinks=true,urlcolor=blue}
\setcounter{secnumdepth}{0}
\setlength{\parindent}{0pt}
\setlength{\parskip}{5pt}
\setlength{\emergencystretch}{3em}
\newcommand{\R}{\mathbb R}
\newcommand{\N}{\mathbb N}
\newcommand{\Z}{\mathbb Z}
\newcommand{\Q}{\mathbb Q}
\newcommand{\C}{\mathbb C}
\newcommand{\D}{\mathbb D}
\newcommand{\F}{\mathbb F}
\newcommand{\sref}[2]{\hyperlink{source:#1:#2}{[S#2]}}
\newcommand{\cataloguescope}[1]{#1}
\begin{document}
\section*{QMA hardness of continuum electronic energy with binary nuclear charges}\label{30007211}
\subsection{Definitions and mathematical statement}
\paragraph{Statement recorded in the supplied source.}
Is the ground-energy promise problem for the full antisymmetric continuum electronic Coulomb Hamiltonian QMA-hard under deterministic classical polynomial-time many-one reductions when nuclear charges are positive binary integers, positions and thresholds are rational, electron count is unary, and the threshold gap is at least inverse input length? Ask for reductions with gap at least one, without auxiliary effective interactions.
\subsection{Short English statement}
Is the ground-energy promise problem for the full antisymmetric continuum electronic Coulomb Hamiltonian QMA-hard under deterministic classical polynomial-time many-one reductions when nuclear charges are positive binary integers, positions and thresholds are rational, electron count is unary, and the threshold gap is at least inverse input length? Ask for reductions with gap at least one, without auxiliary effective interactions.
\subsection{Sources}
\begin{itemize}
\item[S1] Original question source or mathematical context.
\url{https://github.com/openai/math/blob/adc7f1241b42e322a6451854ab7e4b4c146bf78a/preprints/Continuum-Coulomb-hardness-with-binary-nuclear-charges-September-24-2026/build/sections/00-introduction.tex}.
\item[S2] OpenAI, manuscript citation page. Source of the formulation motivating this local problem.
\url{https://github.com/openai/math/blob/adc7f1241b42e322a6451854ab7e4b4c146bf78a/preprints/Continuum-Coulomb-hardness-with-binary-nuclear-charges-September-24-2026/README.md}.
\end{itemize}
\end{document}
